First prove the scale identity.  For \(s>0\), if \(X_k=h_k+\sqrt{v_k/(sr_k)}Z_k\), multiply the observation and parameter by \(\sqrt s\).  The transformed experiment has rates \(r\), and its simple-regret loss is multiplied by \(\sqrt s\).  Taking the infimum over the corresponding one-to-one class of randomized selectors and the supremum over \(H_A\), which is closed under scalar multiplication, gives
\[
 G_A(v_A,sr_A)=s^{-1/2}G_A(v_A,r_A). \tag{18}
\]
For a Bernoulli design, \((Bp)_A=s_A(p)a_A(p)\), so (18) proves the first displayed identity.

Every \(B_{k\ell}\) is positive, hence \(s_A(p)>0\).  Moreover, under nominal label \(\ell\), \(\sum_{k\in A}B_{k\ell}\) is the probability that a binomial count lies in the proper subset \(\{m_k:k\in A\}\) of \(\{0,\ldots,n\}\).  Since \(0<\pi_\ell<1\), counts zero and \(n\) both have positive probability and are absent from the menu.  Therefore
\[
 0<s_A(p)\leq s_{A,\max}<1. \tag{19}
\]
Also \(a_A(p)\in\Delta_A^\circ\).  By the definition of the exact-count value and (18)--(19), for every \(p\),
\[
 G_A(v_A,(Bp)_A)
 =s_A(p)^{-1/2}G_A(v_A,a_A(p))
 \geq s_{A,\max}^{-1/2}R_{CR,A}^\star(v_A). \tag{20}
\]
Taking the infimum over \(p\) proves the weak inequality.  The exact-count value is strictly positive.  Indeed, choose active \(i\ne j\) and restrict nature to the symmetric alternatives \(h=\pm d(e_i-e_j)/2\).  The likelihood-ratio statistic has information \(J=r_i/v_i+r_j/v_j\), so maximizing the equal-prior testing regret over \(d\) gives \(2\kappa/\sqrt J\).  Since \(r_i,r_j\leq1\), this is at least \(2\kappa/(v_i^{-1}+v_j^{-1})^{1/2}>0\).  Thus division by \(R_{CR,A}^\star(v_A)\) is valid.  Combining this fact with \(s_{A,\max}<1\) yields the strict ratio.

For the equality characterization, both inequalities used in (20) must be equalities.  The map \(p\mapsto(Bp)_A\) ranges over a compact set bounded away from zero in every active coordinate.  Gaussian densities vary continuously there, and truncating to a large contrast ball followed by the Gaussian tail bound makes this continuity uniform for the minimax value; hence the Bernoulli infimum is attained.  If the ratio equals the lower bound at an optimizer \(p\), equality in the first inequality of (20) forces \(s_A(p)=s_{A,\max}\), and equality in the second forces \(a_A(p)\) to attain the exact-count infimum.  A convex combination attains the maximum of the linear functional \(s_A\) exactly when it is supported on maximizing columns.  Conversely, any optimizer with those two properties makes both inequalities equalities and attains the ratio.  Inactive coordinates never enter (18)--(20), so they may have zero allocation.